% This is a template for your written document.
%
% To compile using latexmk on the command line, run the following: 
% latexmk -pdf main.tex

\documentclass[12pt]{article}
\usepackage{setspace}
\singlespace
\usepackage[left=1in,right=1in,top=1in,bottom=1in]{geometry}
\usepackage{graphicx}

\title{\textbf{Architectural Tradeoffs in Low-Latency
Algorithmic Trading with Predictive Signals}}
\author{Abhinav Randive}

\begin{document}

\maketitle

\section{Introduction}
Financial markets have increasingly become dominated by automated trading systems that make decisions and execute trades at extremely high speeds. 
These systems rely on algorithmic strategies that analyze market data and react within milliseconds or even microseconds. 
As a result, the design of trading infrastructure has become just as important as the predictive models used to generate trading signals. 
Even when predictive models are accurate, delays in processing or order execution can eliminate any potential advantage in highly competitive markets.
Real-time financial data systems can process very high message volumes, and stream-processing
work in computer systems reports workloads on the order of millions of events per second in
latency-sensitive settings~\cite{stonebrakerLatency}.

Recent advances in machine learning have led to significant interest in applying predictive models to financial markets. 
Researchers have explored techniques ranging from classical statistical models to deep learning architectures for forecasting stock price movements. 
These models often attempt to identify patterns in historical price data, macroeconomic indicators, or market sentiment signals in order to generate short-term predictions about market behavior.

However, predictive accuracy alone does not guarantee profitability in financial markets.
In practice, the ability to act on predictions depends on the latency of the trading system and the structure of the market itself. 
Modern electronic markets operate using limit order books and price-time priority mechanisms, meaning that orders arriving even slightly later than competitors may lose execution priority. 
This creates a critical interaction between prediction, system architecture, and market microstructure.

The goal of this project is to study these architectural tradeoffs by developing an event-driven trading simulation system. 
The system replays historical market data, processes events sequentially, generates short-term predictive signals, and measures the impact of latency and execution constraints on simulated trading outcomes. 
By focusing on the infrastructure that connects predictive models with market execution, this work aims to bridge the gap between machine learning research and practical trading system design.

In addition to predictive modeling challenges, the operational environment of financial markets introduces additional constraints. 
Modern exchanges operate at extremely high speeds and rely on complex technological infrastructures that process millions of messages per second. 
Trading systems must be designed to handle high-frequency streams of constant market data while maintaining a deterministic order and minimizing latency.

These constraints mean that the architecture of the trading system itself plays a crucial role in determining whether predictive signals can actually be translated into profitable trades. 
Even small delays in processing market data or submitting orders can result in lost queue position, reduced fill rates, and increased exposure to adverse price movements. 
Moreover, understanding the interaction between predictive models and trading system architecture has become an important area of study in quantitative finance and financial engineering.

\section{Background}
Understanding algorithmic trading systems requires familiarity with the structure of modern financial markets. 
Most electronic exchanges operate using a limit order book, where buy and sell orders are organized by price and time of submission. 
When traders submit limit orders, they specify a price at which they are willing to buy or sell an asset. 
These orders are placed into a queue and executed when matching orders arrive on the opposite side of the market.

The concept of price-time priority plays a central role in these systems. 
Orders with better prices are executed first, but when multiple orders exist at the same price level, they are executed in the order they were received. 
This means that execution speed can significantly affect trading outcomes. Even small delays in order submission can cause a trader to lose queue position, reducing the likelihood of execution. 
Harris provides a detailed overview of these market microstructure mechanisms and their implications for trading strategies~\cite{harrisTrading}.

HFT trading systems are specifically designed to minimize such delays. 
According to Aldridge, high-frequency trading firms routinely invest heavily in low-latency infrastructure, including using specialized hardware, software, and proximity hosting, trying to be as close to the exchange servers as possible~\cite{aldridgeHFT}. 
These trading systems monitor and process incoming market data and update trading decisions in real time in microseconds. 
The performance of such systems depends not only on the accuracy of prediction models but also on the efficiency of event processing, memory access, and network communication.

Cartea et al.~\cite{carteaMicrostructure} further emphasize that successful algorithmic trading strategies must account for both predictive signaling (BUY/SELL/HOLD) and market impact. 
Trades can influence market prices, and poorly designed strategies may suffer from adverse selection, where informed traders exploit predictable behavior. 
As a result, realistic simulation environments should include the main aspects of market microstructure when evaluating a trading algorithm.

\section{Related Work}

Research on algorithmic trading spans several disciplines, including 
machine learning, financial economics, and computer systems. 
Prior work in this area can broadly be divided into three categories: 
machine learning approaches for market prediction, studies of market 
microstructure, and research on real-time systems architecture.

\subsection{Machine Learning for Financial Prediction}

A large body of research has explored the use of machine learning 
techniques for predicting financial market behavior. Early approaches 
relied primarily on linear regression models and statistical methods, 
but recent work has increasingly focused on deep learning architectures 
capable of capturing complex temporal patterns in financial data.

Fischer and Krauss demonstrate that Long Short-Term Memory (LSTM) 
neural networks can outperform traditional machine learning models 
when predicting stock market returns using historical price data 
~\cite{fischer2018lstm}. Their results highlight the ability of deep 
learning models to capture long-term dependencies in financial time 
series that may not be easily identified by classical statistical 
techniques.

Similarly, Gu, Kelly, and Xiu conduct a comprehensive empirical study 
of machine learning methods applied to asset pricing~\cite{gu2020empirical}. 
Their research evaluates a wide range of algorithms, including neural 
networks and tree-based models, using a large set of financial predictors. 
The authors find that machine learning approaches can improve predictive 
performance relative to traditional linear models, particularly when 
working with high-dimensional datasets.

Despite these advances, predictive modeling in financial markets remains 
challenging due to the non-stationary nature of financial time series. 
Statistical relationships between variables can change over time as market 
conditions evolve, meaning that models which perform well during 
historical backtesting may degrade when deployed in live trading 
environments. Consequently, predictive accuracy alone is often 
insufficient for evaluating the real-world effectiveness of trading 
strategies.

\subsection{Market Microstructure and Trading Dynamics}

Another important area of research focuses on market microstructure, 
which studies the mechanisms through which trades occur in financial 
markets. Modern electronic exchanges operate using limit order books, 
where buy and sell orders are matched according to price-time priority.

Bouchaud, Farmer, and Lillo analyze how markets absorb changes in supply 
and demand over time~\cite{bouchaudTrades}. Their work demonstrates that 
price dynamics emerge from the interaction of orders within the limit 
order book rather than purely from external information shocks. This 
perspective highlights the importance of understanding liquidity and 
order flow when designing algorithmic trading strategies.

Recent research has also explored the integration of alternative data 
sources into predictive models. Patsiarikas et al.~\cite{patsiarikas2025macroeconomic} 
investigate the use of macroeconomic indicators, technical signals, 
and sentiment data for stock market forecasting. Their findings suggest 
that combining multiple sources of information can improve predictive 
performance compared to models that rely solely on historical price data.

However, most studies in this area focus primarily on improving forecast 
accuracy rather than examining the operational constraints involved in 
executing trading strategies within real market environments.

\subsection{Systems Architecture and Low-Latency Processing}

Beyond predictive modeling and market microstructure, the design of 
trading system infrastructure has also been studied within the computer 
systems literature. Modern financial markets generate extremely large 
volumes of real-time data, requiring trading platforms to process 
continuous streams of market events with minimal latency.

Stonebraker et al.~\cite{stonebrakerLatency} describe several key 
requirements for real-time stream processing systems, including 
deterministic event ordering, efficient memory usage, and the ability 
to process high-frequency data streams without introducing significant 
processing delays. These principles are directly applicable to 
algorithmic trading systems.

Event-driven architectures are commonly used in trading systems because 
they allow market activity to be represented as sequences of discrete 
events that can be processed in strict chronological order. Such systems 
must ingest market data, generate trading decisions, and submit orders 
within extremely short time windows.

Despite the importance of these architectural considerations, relatively 
few studies examine the interaction between predictive models and the 
software infrastructure used to execute trading strategies. Many machine 
learning studies assume idealized trading conditions in which predictions 
can be executed instantaneously.

\section{Methodology}

This project develops an event-driven trading simulation to analyze the 
interaction between predictive signals and system latency. The system 
replays historical market data, processes events sequentially, generates 
trading signals, and simulates order execution under realistic constraints.

\subsection{Market Data Acquisition}

Historical market data is obtained using the \texttt{yfinance} API, which 
provides minute-level price and volume data for the S\&P 500 index. The 
data is preprocessed to ensure consistency and usability. Specifically, the
pipeline flattens multi-index columns, selects timestamp/price/volume fields,
converts timestamps into Unix time, and exports a cleaned CSV that is consumed
by the replay engine. This workflow provides reproducible, structured input for
all experiments.

\subsection{Event-Driven Architecture}

Market data is represented as a sequence of discrete events. Each event 
is defined using an \texttt{Event} object containing a timestamp, event 
type, and payload.

Events are processed through a priority queue, ensuring deterministic 
ordering. This prevents look-ahead bias and mirrors real-world trading 
systems where events must be handled sequentially.

\subsection{Deterministic Replay Engine}

The replay engine iterates over historical data and converts each row 
into a \texttt{MARKET\_UPDATE} event. This enables controlled simulation 
of real-time market conditions.

Deterministic replay ensures that results are reproducible and that 
performance metrics are not influenced by future information.

\subsection{Trading Strategy}

The strategy layer combines a rolling-window technical indicator model 
(momentum, volatility, volume, RSI, Bollinger bands, and MACD) with a 
logistic-style mapping to a directional score. Entry and exit rules 
implement stop-loss and take-profit thresholds; new short positions are 
not opened so that the portfolio stays long-only while still allowing 
\texttt{SELL} orders to close long positions.

Although still relatively lightweight, this design separates signal 
computation from decision logic and provides a controlled environment for 
evaluating how system latency affects the ability to act on predictive 
signals.

At each market update, the model computes a directional signal 
strength in the interval \([-1, 1]\). A BUY entry is triggered when the 
signal strength is greater than or equal to the configured entry threshold 
and the portfolio is not already long. A SELL action is triggered to close 
an existing long position when the signal strength is less than or equal to 
the negative threshold, or when stop-loss / take-profit exit conditions are 
met. In the current implementation, the strategy is long-only for new entries 
(i.e., SELL is used to exit longs rather than open new short positions).

\subsection{System Design Tradeoffs}

A central focus of this project is understanding the tradeoffs involved 
in designing low-latency trading systems. While predictive models aim 
to generate accurate signals, the effectiveness of these signals depends 
heavily on how quickly and reliably they can be executed within the system.

One key tradeoff exists between system complexity and latency. More 
advanced trading strategies, such as those incorporating machine learning 
models or multiple data sources, can improve predictive accuracy. However, 
these approaches often require additional computation time, which increases 
latency. In highly competitive markets, even small delays can result in 
loss of execution priority, reducing the likelihood of order fills.

Another important tradeoff involves determinism versus flexibility. 
Event-driven architectures enforce strict ordering of events, ensuring 
that the system processes market data in a consistent and reproducible 
manner. This is essential for accurate backtesting and analysis. However, 
strict ordering can introduce bottlenecks if the system is not optimized 
for high-throughput processing. Designing systems that maintain 
determinism while minimizing delay is therefore a key challenge.

The choice of execution model also introduces tradeoffs. In simplified 
simulations, orders are often assumed to be executed immediately at the 
observed market price. While this assumption simplifies analysis, it 
fails to capture important aspects of real-world trading, such as queue 
position, liquidity constraints, and partial fills. By introducing a 
queue-based matching system with probabilistic fills, this project 
captures a more realistic representation of market behavior. However, 
this added realism comes at the cost of increased system complexity.

Latency measurement itself presents design considerations. Measuring 
latency at a fine-grained level provides valuable insights into system 
performance, but frequent measurement can introduce overhead. The system 
must balance the need for detailed performance data with the goal of 
minimizing measurement-induced delays.

Finally, there is a tradeoff between simulation realism and computational 
efficiency. Highly detailed simulations that incorporate order book depth, 
market impact, and transaction costs provide more accurate insights but 
require significantly more computational resources. In contrast, simpler 
models allow for faster experimentation but may overlook critical dynamics 
of real markets.

By explicitly modeling and evaluating these tradeoffs, this project 
demonstrates that successful algorithmic trading depends not only on 
predictive accuracy but also on careful system design. The interaction 
between prediction, latency, and execution ultimately determines whether 
a trading strategy can be effectively deployed in practice.

\section{Implementation}

This section describes how the conceptual pipeline is realized in software and
how each component contributes to deterministic processing, execution realism,
and measurable system behavior.

\subsection{System Architecture}

The system is composed of modular components that interact through an
event-driven pipeline. The replay engine streams historical data into the
dispatcher, which preserves chronological processing order. The strategy module
produces signals that are transformed into orders by the order manager, then
routed to the execution engine and limit order book for fill simulation under
spread/liquidity constraints. Portfolio and latency-tracking modules record
financial outcomes and performance metrics at each stage. This modular design
allows individual components to be modified and 
evaluated independently. The measured relationship between added strategy
logic and latency overhead is summarized in Fig.~\ref{fig:complexity-latency}.

\begin{figure}[ht]
\centering
\includegraphics[width=0.9\textwidth]{../code/visualizations/04_complexity_tradeoff.png}
\caption{Measured tradeoff between strategy complexity and end-to-end processing latency in the event-driven simulator.}
\label{fig:complexity-latency}
\end{figure}

\subsection{Order Execution Model}

Unlike simplified simulations where orders are executed instantly, this 
system introduces a queue-based execution model.

Orders are placed into buy and sell queues and matched based on price 
conditions and probabilistic fill rules. This models real-world trading 
behavior, where execution depends on queue position and market liquidity.

\subsection{Latency Measurement}

Latency is measured at the event-processing level. For each event, the 
system records stage-level timing for market updates, signal generation, and
order execution, along with end-to-end event latency. This enables quantitative
evaluation of how architectural decisions affect overall system performance and
which stages dominate runtime.

\section{Results}

The system was evaluated using historical S\&P 500 data at one-minute 
intervals. A total of approximately 1,800 events were processed during 
a typical simulation run. Stage-level latency behavior is shown in
Fig.~\ref{fig:latency-impact}.

\begin{figure}[ht]
\centering
\includegraphics[width=0.9\textwidth]{../code/visualizations/03_latency_impact_analysis.png}
\caption{Per-stage latency contribution analysis for the instrumented trading pipeline.}
\label{fig:latency-impact}
\end{figure}

Latency measurements indicate that the system operates efficiently under 
the current configuration. In a representative run, the average processing 
latency per event was 0.0102 milliseconds, with a maximum observed latency 
of 0.0593 milliseconds and p99 latency of 0.0203 milliseconds. These 
results demonstrate that the system can process market data in near real-time 
under the simulated conditions.

From a trading perspective, the strategy generated BUY and SELL decisions 
throughout the run, but only a small number of trades were executed 
(8 fills total). The final portfolio value was \$101,247.09 from an initial 
\$100,000, corresponding to a total return of +1.25\%.

These results indicate that the system is functionally correct and 
operationally efficient, and that profitability is possible in the tested 
configuration. However, because the trade count is low, performance metrics 
should be interpreted cautiously. The return-versus-cost tradeoff across
tested configurations is shown in Fig.~\ref{fig:efficiency-frontier}.

\begin{figure}[ht]
\centering
\includegraphics[width=0.9\textwidth]{../code/visualizations/05_efficiency_frontier.png}
\caption{Efficiency frontier illustrating the relationship between portfolio performance and execution-related costs.}
\label{fig:efficiency-frontier}
\end{figure}

The measured outputs indicate that the current configuration is operationally
stable, with low latency and a modest positive return in the representative
run, while still exhibiting low trade frequency.

\section{Analysis and Discussion}

The results of the simulation provide important insights into the interaction 
between trading strategy design and system architecture.

First, the low average latency confirms that the event-driven architecture 
is effective for processing high-frequency market data streams. The use of 
a priority queue with deterministic ordering ensures that events are handled 
correctly while maintaining performance efficiency.

In this configuration, the strategy produced a small positive return, but 
the low number of fills makes the estimate noisy. This suggests that simple 
technical-indicator rules can be directionally useful, but they remain 
fragile in noisy markets once execution frictions and fill uncertainty are 
included.

The results also highlight the importance of execution modeling. The inclusion 
of probabilistic fills introduces a more realistic representation of market 
conditions, where not all orders are executed immediately. This significantly 
impacts profitability and demonstrates that execution assumptions play a 
critical role in evaluating trading strategies.

To answer the core strategy question explicitly: BUY decisions occur when the 
model confidence exceeds the positive entry threshold, while SELL decisions 
occur when confidence crosses the negative threshold against an open long or 
when risk exits (stop-loss/take-profit) are triggered. This rule-based mapping 
from model output to actions makes the strategy behavior transparent and 
reproducible.

The execution path also applies commission, slippage, and a simple market 
impact estimate on fills. These costs reduce net performance compared with 
a frictionless benchmark and align with the literature on gaps between 
theoretical signals and realized trading outcomes.

A key limitation of the current implementation is the simplicity of the 
strategy and the absence of richer predictive features. Future work could 
explore machine learning-based signals, order book depth analysis, and 
multi-asset strategies to improve performance.

Another limitation is metric robustness under small samples. For example, 
risk-adjusted and execution-quality summaries (such as Sharpe ratio and 
slippage proxies) can become unstable when only a few trades occur, so those 
figures should be treated as diagnostic rather than definitive.

Overall, the findings reinforce the central argument of this project: 
predictive accuracy alone is not sufficient. The ability to execute trades 
efficiently within a low-latency system is equally important in determining 
success in algorithmic trading.

These observations are consistent with prior research in market 
microstructure, which emphasizes that trading outcomes are shaped not 
only by predictive signals but also by execution timing and market 
interaction~\cite{bouchaudTrades, carteaMicrostructure}. The results 
support the argument that bridging the gap between prediction and 
execution is a critical challenge in modern algorithmic trading systems.

\section{Conclusion}
This project presented the design and implementation of an event-driven 
trading simulation framework for analyzing architectural tradeoffs in 
low-latency algorithmic trading systems. By integrating market data replay, 
signal generation, order execution, and latency measurement into a unified 
pipeline, the system provides a controlled environment for studying how 
predictive models interact with execution constraints.

The results demonstrate that while the system achieves low processing 
latency and deterministic event handling, trading performance depends 
critically on execution dynamics and strategy design. In the current test, 
the indicator-driven strategy produced a modest positive return (+1.25\%), 
but with limited trade frequency, highlighting the difficulty of drawing 
strong conclusions about profitability in realistic market environments.

A key contribution of this work is the emphasis on the interaction between 
prediction and execution. Much of the existing literature focuses on improving 
predictive accuracy, often assuming idealized trading conditions. In contrast, 
this project shows that system-level factors such as latency, event ordering, 
and execution modeling play an equally important role in determining outcomes.

Several limitations should be acknowledged. The execution engine uses a 
simplified probabilistic matching model and does not fully capture full order 
book depth, cancellations, or venue-specific fee schedules. The signal model 
remains intentionally lightweight compared with modern machine learning 
approaches or rich multi-factor inputs.

Future work could deepen microstructure realism (explicit price levels, 
cancellations, and calibrated liquidity), expand prediction with trained 
models under latency budgets, and study tail latency under higher event 
throughput.


\bibliographystyle{acm}
\bibliography{bibliography.bib}

\end{document}
